\subsection{拓扑群与带算子空间的逆极限}

我们将说，逆系统 \((\mathbf{G}_{\alpha}, f_{\alpha\beta})\) 是一个拓扑群的逆系统，如果诸 \(G_{\alpha}\) 是拓扑群且诸 \(f_{\alpha\beta}\) 是连续同态。此时 \(G = \varprojlim G_{\alpha}\) 是积群 \(\prod_{\alpha} G_{\alpha}\) 的一个子群；若给 G 赋以由 \(\prod_{\alpha} G_{\alpha}\) 诱导的拓扑群结构，则如此得到的拓扑群称为拓扑群的逆系统 \((\mathbf{G}_{\alpha}, f_{\alpha\beta})\) 的逆极限。若诸 \(G_{\alpha}\) 是豪斯多夫的（相应地，豪斯多夫且完备的），则 G 是豪斯多夫的且在 \(\prod_{\alpha} G_{\alpha}\) 中闭（相应地，豪斯多夫且完备的）（第一章，§ 8，第 2 节，命题 7 的推论 2 以及第二章，§ 3，第 5 节，命题 10 的推论）。

若 \((\mathbf{G}_{\alpha}^{\prime}, f_{\alpha\beta}^{\prime})\) 是另一个拓扑群的逆系统，且对每个 \(\alpha\) ，\(u_{\alpha}: G_{\alpha} \to G_{\alpha}^{\prime}\) 是一个连续同态，使得诸 \(u_{\alpha}\) 构成映射的逆系统，则 \(u = \varprojlim u_{\alpha}\) 是 G 到 \(G^{\prime} = \varprojlim G_{\alpha}^{\prime}\) 中的一个连续同态（第一章，§ 4，第 4 节）。当把“拓扑群”换成“拓扑环”时，同样的结果成立；关于拓扑模的类似结果，我们留给读者去陈述（§ 6，第 6 节）。

设 \((\mathbf{X}_{\alpha}, g_{\alpha \beta})\) 是拓扑空间的逆系统，\((\mathbf{G}_{\alpha}, f_{\alpha \beta})\) 是拓扑群的逆系统。设每个 \(\mathbf{G}_{\alpha}\) 连续地作用于 \(\mathbf{X}_{\alpha}\) （§ 2，第 4 节），且关系 (1) 对所有 \(x_{\beta} \in \mathbf{X}_{\beta}\) 、\(s_{\beta} \in \mathbf{G}_{\beta}\) 与 \(\alpha \leqslant \beta\) 成立。我们已经看到（第 1 节），\(\mathbf{X} = \varprojlim \mathbf{X}_{\alpha}\) 以 \(\mathbf{G} = \varprojlim \mathbf{G}_{\alpha}\) 为算子群；而且 \(\mathbf{G}\) 连续地作用于 \(\mathbf{X}\) 。因为，若 \(g_{\alpha}\) （相应地，\(f_{\alpha}\) ）是典范映射 \(\mathbf{X} \to \mathbf{X}_{\alpha}\) （相应地，\(\mathbf{G} \to \mathbf{G}_{\alpha}\) ），则按定义

\[
g _ {\alpha} (s, x) = f _ {\alpha} (s). g _ {\alpha} (x)
\]

因而诸映射 \((s.x) \to g_{\alpha}(s.x)\) 在 \(\mathbf{X} \times \mathbf{G}\) 上连续，这表明 \((s, x) \to (s.x)\) 的连续性（第一章，§ 4，第 4 节）。

映射 \(h_{\alpha}: \mathbf{X} / \mathbf{G} \to \mathbf{X}_{\alpha} / \mathbf{G}_{\alpha}\) （它由 \(f_{\alpha}\) 与 \(g_{\alpha}\) 诱导）因此是连续的（§ 2，第 4 节），而映射 \(h: \mathbf{X} / \mathbf{G} \to \varprojlim \mathbf{X}_{\alpha} / \mathbf{G}_{\alpha}\) （它由诸 \(h_{\alpha}\) 定义）也是连续的（第一章，§ 2，第 3 节，命题 4）。

命题 1. 设诸 \(\mathbf{X}_{\alpha}\) 与诸 \(\mathbf{G}_{\alpha}\) 满足上述假设。a) 若 \(\mathbf{X}_{\alpha}\) 的每一点的稳定子群都是 \(\mathbf{G}_{\alpha}\) 的一个紧子群（对每个 \(\alpha \in \mathbf{I}\) ），则每一点 \(x = (x_{\alpha})\) （属于 \(\mathbf{X}\) ）的稳定子群是 \(\mathbf{G}\) 的一个紧子群；\(x\) 的轨道（关于 \(\mathbf{G}\) 的）典范地同胚于诸 \(x_{\alpha}\) 的轨道（关于 \(\mathbf{G}_{\alpha}\) 的）的逆极限；且典范映射 \(h: \mathbf{X} / \mathbf{G} \to \varprojlim \mathbf{X}_{\alpha} / \mathbf{G}_{\alpha}\) 是单射的。

\begin{enumerate}
\def\labelenumi{\alph{enumi})}
\setcounter{enumi}{1}
\tightlist
\item
  若对每个 \(\alpha \in \mathbf{I}\) ，\(\mathbf{X}_{\alpha}\) 的点的每个轨道（关于 \(\mathbf{G}_{\alpha}\) 的）都是紧的，则 \(\mathbf{X}\) 的点的每个轨道（关于 \(\mathbf{G}\) 的）都是相对紧的，且 \(h\) 是满射的。若此外 \(h\) 还是双射的，则 \(\mathbf{X}\) 的点的每个轨道都是紧的。
\end{enumerate}

设 \(x = (x_{\alpha}) \in \mathbf{X}\) ，且对每个 \(\alpha \in \mathbf{I}\) ，设 \(\mathbf{X}_{\alpha}' = \mathbf{G}_{\alpha}.x_{\alpha}\) 为 \(x_{\alpha}\) 的轨道。若 \(\alpha \leqslant \beta\) ，则由 (1) 及关系 \(g_{\alpha\beta}(x_{\beta}) = x_{\alpha}\) 推出 \(g_{\alpha\beta}(\mathbf{X}_{\beta}') \subset \mathbf{X}_{\alpha}',\) 因而 \((\mathbf{X}_{\alpha}')\) 是诸 \(\mathbf{X}_{\alpha}\) 的子集的一个逆系统。对每个 \(\alpha \in \mathbf{I}\) ，设 \(u_{\alpha}: \mathbf{G}_{\alpha} \to \mathbf{X}_{\alpha}'\) 是连续映射 \(s_{\alpha} \to s_{\alpha}.x_{\alpha}\) ；诸 \(u_{\alpha}\) 构成映射的逆系统，而 \(u = \varprojlim u_{\alpha}\) 是连续映射 \(s \to s.x\) ，它把 \(G\) 映入子空间 \(X' = \varprojlim X_{\alpha}'\) （含于 \(X\) ）。a) 的假设蕴含 \(\overline{u}_{\alpha}^{-1}(y_{\alpha})\) 对每个 \(y_{\alpha} \in X_{\alpha}'\) 是紧的。由于 \(u_{\alpha}\) 还是满射的，第一章 § 9 第 6 节命题 8 推论 2 的条件得到满足，因而 a) 的前两个断言成立。于是，若 \(x = (x_{\alpha})\) 与 \(y = (y_{\alpha})\) 使得 \(x_{\alpha}\) 与 \(y_{\alpha}\) 属于关于 \(G_{\alpha}\) 的同一轨道（对每个 \(\alpha \in I\) ），则 \(x\) 与 \(y\) 属于关于 \(G\) 的同一轨道，因此 \(h\) 是单射的。

同样，b) 的假设蕴含：典范映射 \(v_{\alpha}: \mathbf{X}_{\alpha} \to \mathbf{X}_{\alpha}/\mathbf{G}_{\alpha}\) 的逆系统满足第一章 § 9 第 6 节命题 8 推论 2 的条件，因而其逆极限 \(v = \varprojlim v_{\alpha}: \mathbf{X} \to \varprojlim \mathbf{X}_{\alpha}/\mathbf{G}_{\alpha}\) 是满射的，且每一点在 \(v\) 下的逆像都是紧的，这里点取遍 \(\varprojlim X_{\alpha}/G_{\alpha}\) 。由于 \(v\) 分解为

\[
\mathrm{X} \xrightarrow {\psi} \mathrm{X} / \mathrm{G} \xrightarrow {h} \varprojlim \mathrm{X} _ {\alpha} / \mathrm{G} _ {\alpha},
\]

其中 \(\psi\) 是典范映射，b) 的各断言由此得出。

推论 1. 若诸 \(\mathbf{G}_{\alpha}\) 是紧的而诸 \(\mathbf{X}_{\alpha}\) 是豪斯多夫的，则 a) 与 b) 的结论成立。

因为 a) 与 b) 的假设都满足：\(G_{\alpha}\) 的每个闭子群都是紧的，且 \(u_{\alpha}: s_{\alpha} \to s_{\alpha}.x_{\alpha}\) 是紧空间到豪斯多夫空间中的连续映射。

推论 2. 若对每个 \(\alpha \in I\) ，群 \(G_{\alpha}\) 传递地作用于空间 \(X_{\alpha}\) ，且 \(X_{\alpha}\) 的每一点的稳定子群都是 \(G_{\alpha}\) 的紧子群，则 \(G\) 传递地作用于 \(X\)，且 \(X\) 的每一点的稳定子群都是 \(G\) 的紧子群。

因为假设 \(a)\) 满足，且 \(\mathbf{X}_{\alpha}^{\prime} = \mathbf{X}_{\alpha}\) （对每个 \(\alpha\) ）。

推论 3. 设诸 \(\mathbf{G}_{\alpha}\) 是豪斯多夫的。对每个 \(\alpha \in \mathbf{I}\) ，设 \(\mathbf{K}_{\alpha}\) 是 \(\mathbf{G}_{\alpha}\) 的一个紧子群，使得 \(f_{\alpha \beta}(\mathbf{K}_{\beta}) \subset \mathbf{K}_{\alpha}\) 对一切 \(\alpha \leqslant \beta\) 成立。那么，若 \(\mathbf{K} = \varprojlim \mathbf{K}_{\alpha}\) ，则典范映射 \(h\) （它把齐性空间 \(\mathbf{G} / \mathbf{K}\) 映到 \(\varprojlim \mathbf{G}_{\alpha} / \mathbf{K}_{\alpha}\) 中）是一个同胚。

\(h\) 是双射这一事实由推论 1 得出，只需把 \(\mathbf{X}_{\alpha}\) 换成 \(\mathbf{G}_{\alpha}\) ，并把 \(\mathbf{G}_{\alpha}\) 换成以右平移作用的 \(\mathbf{K}_{\alpha}\) （§ 2，第 5 节）。设 \(\varphi\) 是典范映射 \(\mathbf{G} \to \mathbf{G} / \mathbf{K}\) ，且对每个 \(\alpha\) ，设 \(f_{\alpha}\) 是典范映射 \(\mathbf{G} \to \mathbf{G}_{\alpha}\) 。若对每个 \(\alpha, \mathbf{V}_{\alpha}\) 取遍 \(e_{\alpha}\) （\(\mathbf{G}_{\alpha}\) 的单位元）的邻域组成的一个基本系，则诸集合 \(\mathbf{V} = \overline{f}_{\alpha}^{-1}(\mathbf{V}_{\alpha}) (\alpha \text{与} \mathbf{V}_{\alpha} \text{变动})\) 构成 \(e\) （\(\mathbf{G}\) 的单位元）的邻域基本系（第一章，§ 4，第 4 节，命题 9），而诸集合 \(\varphi(\mathbf{V}.K)\) 构成 \(\varphi(e)\) 在 \(G/K\) 中的邻域基本系。我们要证明 \(\varphi(V.K)\) 在 \(h\) 下的像包含 \(h(\varphi(e))\) 的一个邻域，亦即存在 \(\beta \geqslant \alpha\) 和一个邻域 \(W_{\beta}\) （\(e_{\beta}\) 在 \(G_{\beta}\) 中），使得 \(\overline{f}_{\beta}^{-1}(W_{\beta}.K_{\beta}) \subset V.K\) 。现在，关系 \(x \in V.K\) 等价于存在 \(y\) （属于 \(K\) ），使 \(f_{\alpha}(xy^{-1}) \in V_{\alpha}\) ，即 \(f_{\alpha}(x) \in V_{\alpha}.f_{\alpha}(K)\) ；从而 \(V.K = \overline{f}_{\alpha}^{-1}(V_{\alpha}.f_{\alpha}(K))\) 。令 \(U_{\alpha} = V_{\alpha}.f_{\alpha}(K)\) ；我们将看到，存在 \(\beta \geqslant \alpha\) ，使得当令 \(U_{\beta} = \overline{f}_{\alpha\beta}^{-1}(U_{\alpha})\) 时有 \(K_{\beta} \subset U_{\beta}\) ；于是将可推出存在一个邻域 \(W_{\beta}\) （\(e_{\beta}\) 在 \(G_{\beta}\) 中），使 \(W_{\beta}.K_{\beta} \subset U_{\beta}\) （第二章，§ 4，第 3 节，命题 4 的推论），而这便建立了所要的关系

\[
\overline {{f}} _ {\beta} ^ {-1} (\mathrm{W} _ {\beta}. \mathrm{K} _ {\beta}) \subset \overline {{f}} _ {\beta} ^ {-1} (\mathrm{U} _ {\beta}) = \mathrm{V}. \mathrm{K}.
\]

我们用反证法来进行：对每个 \(\beta \geqslant \alpha\) ，用 \(\mathbf{M}_{\beta}\) 表示 \(\mathbf{K}_{\beta} \cap \mathbb{C}\mathbf{U}_{\beta}\) ；由于 \(\overline{f}_{\beta \gamma}^{-1}(\mathrm{U}_{\beta}) = \mathrm{U}_{\gamma}\) （当 \(\alpha \leqslant \beta \leqslant \gamma\) 时），诸 \(\mathbf{M}_{\beta}\) 构成诸 \(\mathbf{G}_{\beta}\) 的紧子集的一个逆系统（对 \(\beta \geqslant \alpha\) ）。如果它们都非空，则它们的逆极限 \(\mathbf{M}\) 也非空（第一章，§ 9，第 6 节，命题 8）。显然 \(\mathbf{M} \subset \mathbf{K}\) 且 \(f_{\alpha}(\mathbf{M}) \subset \mathbf{M}_{\alpha}\) ；但这是荒谬的，因为 \(f_{\alpha}(\mathbf{K}) \subset \mathbf{U}_{\alpha}\) ，于是证明完毕。

